\section{随机森林填补}

\begin{definition}[随机森林插值]
    具体做法是为每个缺失值特征训练一个随机森林模型，通过该模型对缺失值进行预测。
\end{definition}

\begin{tips}
    \begin{itemize}
        \item 优点：能够处理非线性关系的特征，预测准确性较高。
        \item 缺点：需要额外的训练过程，可能增加计算复杂度。
    \end{itemize}
\end{tips}

\subsection{数据预处理阶段的随机森林填充}

此时我们的目标主要是修复数据，此时没有 \texttt{Y\_full} 这个值，我们的目的是修复数据。

\subsection{模型训练阶段}

如果我们拥有一个特征表 X 和目标 Y，我们的目标是预测 Y，此时 Y 是标准答案。
